% math0 manuscript. Build with: latexmk -pdf paper.tex
\documentclass[11pt]{amsart}
\usepackage[margin=1.1in]{geometry}
\usepackage{amsmath,amssymb,amsthm,mathtools}
\usepackage{url}
\usepackage{hyperref}
\usepackage{xcolor}
\usepackage{listings}
\usepackage{newunicodechar}

\lstdefinelanguage{lean}{
  keywords={theorem,lemma,def,abbrev,example,fun,by,have,show,from,let,in,if,then,else,match,with,
    namespace,end,open,import,section,variable,instance,structure,where,calc,obtain,intro,exact,
    refine,rcases,rintro,simp,omega,decide,norm_num,linarith,nlinarith,positivity,induction,cases,
    private,rw},
  sensitive=true, comment=[l]{--}, morecomment=[s]{/-}{-/}, string=[b]"}
\lstset{language=lean, basicstyle=\ttfamily\footnotesize, columns=fullflexible, keepspaces=true, breaklines=true,
  keywordstyle=\color{blue!60!black}\bfseries, commentstyle=\color{gray}, frame=single,
  extendedchars=true, inputencoding=utf8,
  literate={ℕ}{{$\mathbb{N}$}}1 {ℤ}{{$\mathbb{Z}$}}1 {∀}{{$\forall$}}1 {↔}{{$\leftrightarrow$}}1
    {∧}{{$\wedge$}}1 {→}{{$\to$}}1 {≤}{{$\le$}}1 {≠}{{$\ne$}}1 {∑}{{$\sum$}}1 {∈}{{$\in$}}1
    {×}{{$\times$}}1 {ˢ}{{$^{\mathrm{s}}$}}1}
% Unicode commonly found in Lean sources.
\newunicodechar{ℕ}{\ensuremath{\mathbb{N}}}
\newunicodechar{ℤ}{\ensuremath{\mathbb{Z}}}
\newunicodechar{ℚ}{\ensuremath{\mathbb{Q}}}
\newunicodechar{ℝ}{\ensuremath{\mathbb{R}}}
\newunicodechar{ℂ}{\ensuremath{\mathbb{C}}}
\newunicodechar{∀}{\ensuremath{\forall}}
\newunicodechar{∃}{\ensuremath{\exists}}
\newunicodechar{→}{\ensuremath{\to}}
\newunicodechar{←}{\ensuremath{\leftarrow}}
\newunicodechar{↔}{\ensuremath{\leftrightarrow}}
\newunicodechar{↦}{\ensuremath{\mapsto}}
\newunicodechar{¬}{\ensuremath{\neg}}
\newunicodechar{∧}{\ensuremath{\wedge}}
\newunicodechar{∨}{\ensuremath{\vee}}
\newunicodechar{≤}{\ensuremath{\le}}
\newunicodechar{≥}{\ensuremath{\ge}}
\newunicodechar{≠}{\ensuremath{\ne}}
\newunicodechar{∈}{\ensuremath{\in}}
\newunicodechar{∉}{\ensuremath{\notin}}
\newunicodechar{⊆}{\ensuremath{\subseteq}}
\newunicodechar{∩}{\ensuremath{\cap}}
\newunicodechar{∪}{\ensuremath{\cup}}
\newunicodechar{∑}{\ensuremath{\sum}}
\newunicodechar{∏}{\ensuremath{\prod}}
\newunicodechar{⟨}{\ensuremath{\langle}}
\newunicodechar{⟩}{\ensuremath{\rangle}}
\newunicodechar{·}{\ensuremath{\cdot}}
\newunicodechar{∘}{\ensuremath{\circ}}
\newunicodechar{⁻}{\ensuremath{^{-}}}
\newunicodechar{¹}{\ensuremath{^{1}}}
\newunicodechar{²}{\ensuremath{^{2}}}
\newunicodechar{₀}{\ensuremath{_{0}}}
\newunicodechar{₁}{\ensuremath{_{1}}}
\newunicodechar{₂}{\ensuremath{_{2}}}
\newunicodechar{α}{\ensuremath{\alpha}}
\newunicodechar{β}{\ensuremath{\beta}}
\newunicodechar{ε}{\ensuremath{\varepsilon}}
\newunicodechar{δ}{\ensuremath{\delta}}
\newunicodechar{π}{\ensuremath{\pi}}
\newunicodechar{σ}{\ensuremath{\sigma}}
\newunicodechar{φ}{\ensuremath{\varphi}}
\newunicodechar{∞}{\ensuremath{\infty}}
\newunicodechar{⌊}{\ensuremath{\lfloor}}
\newunicodechar{⌋}{\ensuremath{\rfloor}}
\newunicodechar{⌈}{\ensuremath{\lceil}}
\newunicodechar{⌉}{\ensuremath{\rceil}}
\newunicodechar{∣}{\ensuremath{\mid}}
\newunicodechar{≡}{\ensuremath{\equiv}}
\newunicodechar{×}{\ensuremath{\times}}
\newunicodechar{⧸}{\ensuremath{/}}
\newunicodechar{‖}{\ensuremath{\|}}
\newunicodechar{√}{\ensuremath{\surd}}
\newunicodechar{∅}{\ensuremath{\emptyset}}
\newunicodechar{⦃}{\ensuremath{\{\!|}}
\newunicodechar{⦄}{\ensuremath{|\!\}}}
\newunicodechar{𝓝}{\ensuremath{\mathcal{N}}}
\newunicodechar{≫}{\ensuremath{\gg}}
\newunicodechar{≪}{\ensuremath{\ll}}

\theoremstyle{plain}
\newtheorem{theorem}{Theorem}
\newtheorem{lemma}[theorem]{Lemma}
\newtheorem{proposition}[theorem]{Proposition}
\newtheorem{corollary}[theorem]{Corollary}
\theoremstyle{definition}
\newtheorem{definition}[theorem]{Definition}
\newtheorem{remark}[theorem]{Remark}

\emergencystretch=3em
% Lean identifiers in running text (breakable at dots and underscores).
\DeclareUrlCommand\decl{\urlstyle{tt}}
\expandafter\def\expandafter\UrlBreaks\expandafter{\UrlBreaks\do\_\do\.}
\newcommand{\N}{\mathbb{N}}
\newcommand{\Z}{\mathbb{Z}}
\newcommand{\nr}{\approx}
\newcommand{\monus}{\mathbin{\dot{-}}}
\newcommand{\ind}[1]{[\,#1\,]}

\title[Hardin's formula for OEIS A195249]{Lower triangles of a $3\times3$ array whose king-neighbours differ by at most two: proof of Hardin's formula $a(n)=1047n-1459$ (OEIS A195249)}
\author{math0}
\address{Produced by Claude (an AI model developed by Anthropic) in the math0 project}
\date{October 20, 2026}

\begin{document}
\begin{abstract}
Let $a(n)$ be the number of ways to fill the six cells of the lower triangle of a $3\times3$ array with integers in $\{0,\dots,n\}$ so that any two cells that are horizontal, vertical or diagonal neighbours receive entries differing by at most $2$; this is OEIS A195249.
We prove the formula $a(n)=1047n-1459$ for all $n>2$, which R.~H. Hardin recorded as empirical in 2011; the generating function and the recurrence conjectured by C.~Barker in 2018 follow at once.
The proof splits the count by the entry at the apex; since every entry lies within distance $4$ of the apex, the count for a fixed apex is invariant under raising the upper bound and under translation once the apex is far from both ends of the range, which reduces the formula to a finite computation.
To our knowledge this is the first proof (the OEIS entry still lists the formula as empirical), but the argument is routine: it is essentially the bounded-span argument given by R.~Israel in 2017 for a sibling table.
The theorem is verified in Lean~4 against a statement that math0 formalized from the OEIS entry and froze after a fidelity review, with definitions checked by the kernel against all $38$ listed terms; it is not part of the formal-conjectures repository.
\end{abstract}
\maketitle

\section{Introduction}

The sequence A195249 of the OEIS \cite{A195249} was contributed by R.~H. Hardin on September~13, 2011. Its offset is $1$, its first terms are $64, 729, 1682, 2729, 3776, 4823, 5870,\dots$, and its name reads
\begin{quote}
\small Number of lower triangles of a 3 X 3 0..n array with each element differing from all of its diagonal, vertical, antidiagonal and horizontal neighbors by two or less.
\end{quote}
It is row $3$ of Hardin's table A195248 \cite{A195248}, whose entry in row $m$ and column $k$ counts the lower triangles of an $m\times m$ array with entries in $0..k$ under the same condition (the OEIS quotations below call this entry \texttt{T(n,k)}; it is unrelated to the function $T(n,c)$ used in our proof).
Hardin recorded in A195249 the ``empirical'' formula $a(n)=1047n-1459$ for $n>2$, and in A195248 empirical linear formulas for rows $1$ to $7$ together with the remark \texttt{Generalizing, T(n,k) = A195214(n)*k + const(n) for k>2*n-4}. Here A195214 \cite{A195214} counts the same triangles with entries in $\Z$, two triangles that differ by a constant being counted once; A195214$(3)=1047$.
On May~6, 2018, C.~Barker added to A195249 a conjectured generating function and a conjectured recurrence (both quoted in Section~\ref{sec:statement}).
We prove Hardin's formula (Theorem~\ref{thm:main}) and deduce Barker's two conjectures (Corollary~\ref{cor:barker}).
Let $T(n,c)$ be the number of valid triangles with entries in $0..n$ and apex entry $c$. All their entries lie in $[c-4,c+4]$, so $T(n+1,c)=T(n,c)$ if $c+4\le n$, and $T(n+1,c+1)=T(n,c)$ if $c\ge4$ (Lemmas~\ref{lem:up}, \ref{lem:shift}). Hence $a(n)=a(n-1)+T(8,4)$ for $n\ge8$, and the computations $T(8,4)=1047$ and $a(3),\dots,a(7)$ complete the proof.

\subsection*{Status and prior work}
The closest earlier work is a comment of Robert Israel, dated November~21, 2017, in the sibling table A195232 \cite{A195232}, where neighbouring entries differ by one or less: \texttt{Since elements of a solution differ by no more than n, T(n,k)-T(n,k-1) is constant for k >= n. This confirms the empirical formula: T(n,k) is a polynomial of degree 1 in k for k > n-3.}
Our proof applies the same bounded-span idea to row $3$ of A195248, with explicit thresholds and a finite computation.
In A195220 \cite{A195220}, which counts the triangles with entries in $\Z$, up to constants, when neighbours differ by $k$ or less (its column $2$ is A195214), Israel noted on October~6, 2019, that the count is an Ehrhart polynomial in $k$.
Dougherty-Bliss and Kauers \cite{DBK} proved several of Hardin's OEIS conjectures, among them formulas linear in one dimension of the array, for a different family of king-move arrays that Hardin contributed in 2014; their paper does not mention A195248 or A195249.
In the copy of the OEIS that we used (A195249 at revision \#9 of May~6, 2018), the entries A195249, A195248 and A195214 have no edit notes; the first two still mark their formulas as empirical or conjectural, as does the live entry A195249 as of October 2026, and A195214 contains no formula. The math0 referees found no proof in searches of the web, arXiv and MathOverflow for the formula, the A-numbers and the name of the entry.
The conjecture is not among the OEIS statements of the formal-conjectures repository \cite{FC}.
We therefore classify the result as \emph{new}, with the caveat that it is not new mathematics: given Israel's comment, what remains is to fix the thresholds and to carry out the finite computation, and the contribution of this note is a complete, machine-checked proof.

\section{Statement}\label{sec:statement}

\begin{definition}\label{def:main}
Number the six cells $(i,j)$, $0\le j\le i\le2$ (row $i$, column $j$), of the lower triangle of a $3\times3$ array as
\[
  \begin{aligned}&c_0=(0,0),\quad c_1=(1,0),\quad c_2=(1,1),\\ &c_3=(2,0),\quad c_4=(2,1),\quad c_5=(2,2),\end{aligned}
  \qquad\text{with entries}\qquad
  \begin{array}{ccc} x_0 & & \\ x_1 & x_2 & \\ x_3 & x_4 & x_5 \end{array}
\]
for a filling $x=(x_0,\dots,x_5)\in\N^6$. Two distinct cells are \emph{neighbours} if their rows and their columns each differ by at most $1$ (horizontal, vertical, diagonal and antidiagonal neighbours). The neighbouring pairs $\{k,l\}$ of indices are the ten pairs
\begin{equation}\label{eq:edges}
  01,\ 02,\ 12,\ 13,\ 23,\ 14,\ 24,\ 34,\ 25,\ 45 .
\end{equation}
For $x,y\in\N$ write $x\nr y$ if $|x-y|\le2$. A filling is \emph{valid} if $x_k\nr x_l$ for every neighbouring pair $\{k,l\}$, and $a(n)$ is the number of valid fillings in $\{0,\dots,n\}^6$.
\end{definition}

The FORMULA section of A195249 \cite{A195249} contains, verbatim, Hardin's formula and Barker's conjectures:
\begin{quote}\small\ttfamily
Empirical: a(n) = 1047*n - 1459 for n>2.\\
Conjectures from \_Colin Barker\_, May 06 2018: (Start)\\
G.f.: x*(64 + 601*x + 288*x\^{}2 + 94*x\^{}3) / (1 - x)\^{}2.\\
a(n) = 2*a(n-1) - a(n-2) for n>4.\\
(End)
\end{quote}

\begin{theorem}\label{thm:main}
For every integer $n\ge3$, $a(n)=1047n-1459$.
\end{theorem}

The bound is sharp: $a(2)=3^6=729$, while $1047\cdot2-1459=635$.

\subsection*{Lean formalization}
The statement was formalized by math0 in \path{lean/Math0/Statements/OEIS/A195249.lean}, reviewed for fidelity and frozen (SHA-256 \texttt{ee654c51\ldots fe44c}) before any proof attempt. Its definitions and the statement read (verbatim; the header, the import, comments and docstrings, blank lines, two decidability instances, all theorems, the auxiliary definitions \decl{win}, \decl{lsum} and \decl{aL} (shown in Section~\ref{sec:proof}) and the two propositions recording Barker's conjectures are omitted):
\begin{lstlisting}[basicstyle=\ttfamily\scriptsize]
namespace Math0.Statements.A195249
open Finset
def cell : Fin 6 → ℕ × ℕ := ![(0, 0), (1, 0), (1, 1), (2, 0), (2, 1), (2, 2)]
def Adj (p q : ℕ × ℕ) : Prop :=
  p ≠ q ∧ p.1 ≤ q.1 + 1 ∧ q.1 ≤ p.1 + 1 ∧ p.2 ≤ q.2 + 1 ∧ q.2 ≤ p.2 + 1
def near (x y : ℕ) : Prop := x ≤ y + 2 ∧ y ≤ x + 2
def Valid (v : Fin 6 → ℕ) : Prop := ∀ i j : Fin 6, Adj (cell i) (cell j) → near (v i) (v j)
def a (n : ℕ) : ℕ :=
  ∑ x0 ∈ range (n + 1), ∑ x1 ∈ range (n + 1), ∑ x2 ∈ range (n + 1), ∑ x3 ∈ range (n + 1),
    ∑ x4 ∈ range (n + 1), ∑ x5 ∈ range (n + 1), if Valid ![x0, x1, x2, x3, x4, x5] then 1 else 0
def statement : Prop := ∀ n : ℕ, 2 < n → (a n : ℤ) = 1047 * n - 1459
end Math0.Statements.A195249
\end{lstlisting}
This encodes Definition~\ref{def:main}: \decl{cell} lists $c_0,\dots,c_5$, \decl{Adj} is the neighbour relation and \decl{near} is $\nr$ (both written without truncated subtraction), and \texttt{a n} is the six-fold sum over $\{0,\dots,n\}^6$ of the indicator of validity, i.e.\ the OEIS definition itself. The statement is cast to $\Z$, so $1047n-1459$ is not a truncated difference; it has exactly the hypothesis $n>2$ of the OEIS formula, and as the offset is $1$, the value $a(0)=1$ plays no role. Barker's conjectures are recorded in the same file as separate propositions (\decl{statementBarkerGF}, \decl{statementBarkerRec}), not proved in Lean here.
The file also proves the following sanity checks: \decl{cell_injective} and \decl{image_cell} (by \texttt{decide}) state that \decl{cell} lists each position $(i,j)$ with $j\le i<3$ exactly once; \decl{valid_iff} states that \texttt{Valid ![x0, x1, x2, x3, x4, x5]} is equivalent to the conjunction of the ten conditions \texttt{near xk xl} for the pairs \eqref{eq:edges}; \decl{a_small} states \texttt{a 1 = 64 ∧ a 2 = 729}, evaluated in the kernel straight from the definition (\texttt{decide +kernel}); and \decl{a_data} states that \texttt{(List.range 38).map (fun i => a (i + 1))} is the list of all $38$ terms $64,729,\dots,38327$ listed in the OEIS entry, evaluated in the kernel through the pruned copy \decl{aL} of Section~\ref{sec:proof} and the lemma \texttt{a\_eq\_aL (n : ℕ) : a n = aL n}.

\section{Proof}\label{sec:proof}

The proof follows \path{lean/Math0/Frontier/OEIS/A195249.lean}; we give the Lean name of each step.

\subsection{Windows and the apex decomposition}
For $n,y\in\N$ let
\[
  W_n(y)=\{x\in\N : y\monus2\le x\le\min(y+2,n)\},\qquad y\monus 2=\max(y-2,0).
\]
Then $W_n(y)=\{x\in\{0,\dots,n\}: x\nr y\}$ (\decl{filter_near}). In Lean the window is the list \texttt{win n y} of consecutive integers from $y\monus 2$ to $\min(y+2,n)$, sums over lists are written with \decl{lsum}, and the statement file defines the pruned copy \decl{aL} of \decl{a}:
\begin{lstlisting}
def win (n y : ℕ) : List ℕ := List.range' (y - 2) (min (y + 2) n + 1 - (y - 2))
def lsum (l : List ℕ) (f : ℕ → ℕ) : ℕ := l.foldr (fun x acc => f x + acc) 0
def aL (n : ℕ) : ℕ :=
  lsum (List.range (n + 1)) fun x0 =>
  lsum (win n x0) fun x1 =>
  lsum (win n x0) fun x2 => if near x1 x2 then
  lsum (win n x1) fun x3 => if near x2 x3 then
  lsum (win n x1) fun x4 => if near x2 x4 ∧ near x3 x4 then
  lsum (win n x2) fun x5 => if near x4 x5 then 1 else 0
  else 0 else 0 else 0
\end{lstlisting}
For $n,c\in\N$ define, with $\ind{P}\in\{0,1\}$ the indicator of $P$,
\begin{multline}\label{eq:T}
  T(n,c)=\sum_{x_1\in W_n(c)}\ \sum_{x_2\in W_n(c)}\ind{x_1\nr x_2}\sum_{x_3\in W_n(x_1)}\ind{x_2\nr x_3}\\
  \sum_{x_4\in W_n(x_1)}\ind{x_2\nr x_4,\ x_3\nr x_4}\sum_{x_5\in W_n(x_2)}\ind{x_4\nr x_5}.
\end{multline}
In Lean, \texttt{T n c} is literally the inner function of \decl{aL}.

\begin{lemma}[\decl{a_eq_aL}, \decl{a_eq_sum}]\label{lem:sum}
For every $n$, $a(n)=\sum_{c=0}^{n}T(n,c)$, and $T(n,c)$ is the number of valid fillings $x\in\{0,\dots,n\}^6$ with $x_0=c$.
\end{lemma}

\begin{proof}
By \decl{valid_iff}, validity is the conjunction of the ten conditions $x_k\nr x_l$ for the pairs \eqref{eq:edges}. Attach each condition to the larger index of its pair, so that $x_1,\dots,x_5$ carry $\{01\}$, $\{02,12\}$, $\{13,23\}$, $\{14,24,34\}$, $\{25,45\}$. In the six-fold sum defining $a(n)$, move each condition outwards to the sum over the variable carrying it (\decl{sum_pull}), and use the condition of $x_1,\dots,x_5$ with $x_0,x_0,x_1,x_1,x_2$ respectively to restrict the range $\{0,\dots,n\}$ of that variable to a window (\decl{sum_win}, \decl{sum_win_and}); the other conditions stay as indicators. This gives \eqref{eq:T} for the sum over $x_1,\dots,x_5$ with $x_0=c$ fixed, and is \decl{a_eq_aL} of the statement file; \decl{a_eq_sum} rewrites the outer list sum as $\sum_{c=0}^n$ (\decl{sum_range_eq_lsum}).
\end{proof}

\begin{lemma}\label{lem:win}
Let $n,y,x\in\N$.
\begin{enumerate}
\item[(a)] \textup{(\decl{mem_win})} If $x\in W_n(y)$ then $y\monus2\le x\le y+2$ and $x\le n$.
\item[(b)] \textup{(\decl{win_succ_up})} If $y+2\le n$ then $W_{n+1}(y)=W_n(y)$.
\item[(c)] \textup{(\decl{win_succ_shift})} If $y\ge2$ then $W_{n+1}(y+1)=W_n(y)+1$, where $W_n(y)+1\coloneqq\{x+1: x\in W_n(y)\}$.
\item[(d)] \textup{(\decl{near_succ})} $x+1\nr y+1$ if and only if $x\nr y$.
\end{enumerate}
\end{lemma}

\begin{proof}
(a) is the definition. For (b), $\min(y+2,n+1)=y+2=\min(y+2,n)$. For (c), $y\ge2$ gives $(y+1)\monus2=(y\monus2)+1$, and $\min(y+3,n+1)=\min(y+2,n)+1$; so both windows are intervals, the first being the second shifted by $1$. (d) is clear. In Lean, (b) and (c) are equalities of lists (for (c), \decl{List.map_add_range'}), so sums over the windows can be rewritten with them.
\end{proof}

\subsection{Two stability properties}

\begin{lemma}[upper stability, \decl{T_up}]\label{lem:up}
If $c+4\le n$ then $T(n+1,c)=T(n,c)$.
\end{lemma}

\begin{proof}
Since $c+2\le n$, Lemma~\ref{lem:win}(b) gives $W_{n+1}(c)=W_n(c)$, so in \eqref{eq:T} for $T(n+1,c)$ the outer two sums run over the same windows as for $T(n,c)$. For $x_1,x_2\in W_n(c)$, Lemma~\ref{lem:win}(a) gives $x_1,x_2\le c+2$, hence $x_1+2\le n$ and $x_2+2\le n$, and Lemma~\ref{lem:win}(b) gives $W_{n+1}(x_1)=W_n(x_1)$ and $W_{n+1}(x_2)=W_n(x_2)$. Thus the inner sums also agree, term by term (\decl{lsum_congr}).
\end{proof}

\begin{lemma}[translation, \decl{T_shift}]\label{lem:shift}
If $c\ge4$ then $T(n+1,c+1)=T(n,c)$.
\end{lemma}

\begin{proof}
Since $c\ge2$, Lemma~\ref{lem:win}(c) gives $W_{n+1}(c+1)=W_n(c)+1$, so in \eqref{eq:T} for $T(n+1,c+1)$ we may substitute $x_1=x_1'+1$ and $x_2=x_2'+1$ with $x_1',x_2'\in W_n(c)$ (\decl{lsum_map_succ}). By Lemma~\ref{lem:win}(a), $x_1',x_2'\ge c-2\ge2$, so Lemma~\ref{lem:win}(c) gives $W_{n+1}(x_1'+1)=W_n(x_1')+1$ and $W_{n+1}(x_2'+1)=W_n(x_2')+1$, and we substitute $x_k=x_k'+1$ for $k=3,4,5$ as well. By Lemma~\ref{lem:win}(d), every indicator $\ind{x_k'+1\nr x_l'+1}$ equals $\ind{x_k'\nr x_l'}$, and the resulting expression is \eqref{eq:T} for $T(n,c)$.
\end{proof}

Combinatorially (Lemma~\ref{lem:sum}), Lemma~\ref{lem:up} says that no valid filling with apex $c$ has an entry above $c+4$, and Lemma~\ref{lem:shift} that adding $1$ to every entry maps the valid fillings in $\{0,\dots,n\}^6$ with apex $c$ bijectively onto those in $\{0,\dots,n+1\}^6$ with apex $c+1$.

\subsection{The finite computations}

\begin{lemma}[\decl{T_eight}, \decl{T_four}, \decl{a_vals}]\label{lem:comp}
$T(8,4)=1047$, and hence $T(m+8,4)=1047$ for every $m\ge0$. Moreover $a(3)=1682$, $a(4)=2729$, $a(5)=3776$, $a(6)=4823$ and $a(7)=5870$.
\end{lemma}

\begin{proof}
The values $T(8,4)$ and $a(3),\dots,a(7)$ (the latter as $\mathrm{aL}(3),\dots,\mathrm{aL}(7)$, by \decl{a_eq_aL}) are computed by \texttt{decide +kernel}. The second claim follows by induction on $m$ from Lemma~\ref{lem:up}, which applies since $4+4\le m+8$.
\end{proof}

No certificate is involved: \texttt{decide +kernel} turns the \texttt{Decidable} instance into a proof that the kernel checks by reducing, e.g., \texttt{decide (T 8 4 = 1047)} to \texttt{true}, evaluating the nested sums \eqref{eq:T} with its built-in arithmetic on natural-number literals. Each window has at most $5$ elements, so $T(n,c)$ has at most $5^5=3125$ innermost terms. Unlike \texttt{native\_decide}, this trusts no compiled code and adds no axiom.

\subsection{Constant first differences and the proof of Theorem~\ref{thm:main}}

\begin{lemma}[\decl{step}]\label{lem:step}
For every $m\ge0$, $a(m+8)=a(m+7)+1047$.
\end{lemma}

\begin{proof}
Put $N=m+8$. By Lemma~\ref{lem:sum},
\[
  a(N)=\sum_{c=0}^{3}T(N,c)+T(N,4)+\sum_{c=5}^{N}T(N,c),\qquad
  a(N-1)=\sum_{c=0}^{3}T(N-1,c)+\sum_{c=4}^{N-1}T(N-1,c)
\]
(\decl{split4}, \decl{Finset.sum_range_succ'}). For $c\le3$ we have $c+4\le7\le N-1$, so $T(N,c)=T(N-1,c)$ by Lemma~\ref{lem:up}. For $5\le c\le N$, Lemma~\ref{lem:shift} applied to $c-1\ge4$ gives $T(N,c)=T(N-1,c-1)$, so $\sum_{c=5}^{N}T(N,c)=\sum_{c=4}^{N-1}T(N-1,c)$. Finally $T(N,4)=1047$ by Lemma~\ref{lem:comp}. Hence $a(N)=a(N-1)+1047$.
\end{proof}

\begin{proof}[Proof of Theorem~\ref{thm:main}] (\decl{formula}, \decl{result})
For $3\le n\le7$ the values of Lemma~\ref{lem:comp} agree with $1047n-1459$, which takes the values $1682, 2729, 3776, 4823, 5870$. For $n\ge7$ write $n=m+7$ and induct on $m$: the case $m=0$ is $a(7)=5870=1047\cdot7-1459$, and if $a(m+7)=1047(m+7)-1459$ then Lemma~\ref{lem:step} gives $a(m+8)=1047(m+8)-1459$. In Lean this induction is carried out in $\Z$.
\end{proof}

\subsection{Consequences and a remark (not formalized)}

\begin{corollary}[Barker's conjectures]\label{cor:barker}
For $n>4$, $a(n)=2a(n-1)-a(n-2)$, and
\[
  \sum_{n\ge1}a(n)x^n=\frac{x(64+601x+288x^2+94x^3)}{(1-x)^2}.
\]
\end{corollary}

\begin{proof}
For $n\ge5$ the indices $n,n-1,n-2$ are at least $3$, so the recurrence holds because $a$ is affine there by Theorem~\ref{thm:main}. Multiplying the power series by $(1-x)^2$, the coefficient of $x^k$ is $a(k)-2a(k-1)+a(k-2)$ (with $a(0)$ and $a(-1)$ read as $0$), which is $64$, $729-128=601$, $1682-1458+64=288$, $2729-3364+729=94$ for $k=1,2,3,4$, and $0$ for $k\ge5$ by the recurrence.
\end{proof}

The recurrence fails at $n=4$ ($2a(3)-a(2)=2635$). Corollary~\ref{cor:barker} is not formalized.

\begin{remark}[shapes]\label{rem:shapes}
Call a valid filling with entries in $\Z$ and minimum $0$ a \emph{shape}, and its maximum its \emph{span}. Any two cells are joined by a chain of at most two neighbouring pairs, so every span is at most $4$, and a shape of span $s$ has $\max(0,n+1-s)$ translates in $\{0,\dots,n\}^6$. Hence $a(n)=\sum_{s=0}^{4}N_s\max(0,n+1-s)$, where $N_s$ is the number of shapes of span $s$. A computer enumeration (not part of the Lean development) gives $(N_0,\dots,N_4)=(1,62,602,288,94)$: there are $1047=\text{A195214}(3)$ shapes, of total span $2506$. For $n\ge3$ this gives again $a(n)=1047(n+1)-2506=1047n-1459$ (Israel's bounded-span argument; the slope is A195214$(3)$, as in Hardin's ``Generalizing'' remark); it also gives $\sum_{n\ge1}a(n)x^n=(1+62x+602x^2+288x^3+94x^4)/(1-x)^2-1$, Barker's generating function.
\end{remark}

\section{Machine verification}\label{sec:lean}

\subsection*{Files and declarations}
The proof is the file \path{lean/Math0/Frontier/OEIS/A195249.lean} of the math0 repository, \url{https://github.com/Foundation-0/math0}. It imports the frozen statement file (Section~\ref{sec:statement}) and the audit commands in \path{lean/Math0/Audit.lean}. The main declaration is \decl{Math0.Frontier.A195249.result}, of type \decl{Math0.Statements.A195249.statement}; all other declarations are \texttt{private}, among them the function \decl{T} (the inner function of \decl{aL}) and the lemmas \decl{a_eq_sum}, \decl{T_up}, \decl{T_shift}, \decl{T_eight}, \decl{T_four}, \decl{step}, \decl{a_vals} and \decl{formula}, whose statements are those of Section~\ref{sec:proof}. The proof of \decl{a_vals} is \texttt{simp only [a\_eq\_aL]} followed by \texttt{decide +kernel}, and neither file contains \texttt{sorry}, \texttt{native\_decide}, \texttt{axiom} or \texttt{set\_option}. Verbatim excerpts:
\begin{lstlisting}
private theorem T_eight : T 8 4 = 1047 := by
  decide +kernel
theorem result : statement := ...
#math0_settles Math0.Frontier.A195249.result Math0.Statements.A195249.statement
\end{lstlisting}

\subsection*{Audit}
The command \texttt{\#math0\_settles} (in \path{lean/Math0/Audit.lean}) fails unless (1) \decl{statement} is a \texttt{Prop}-valued definition without universe parameters; (2) \decl{statement} and \decl{result} depend on no axioms besides \decl{propext}, \decl{Classical.choice} and \decl{Quot.sound} (rejecting \decl{sorryAx}, \decl{Lean.ofReduceBool} from \texttt{native\_decide}, and custom axioms); and (3) the type of \decl{result} is syntactically \decl{statement} or its negation, up to metadata. It reports
\begin{lstlisting}[basicstyle=\ttfamily\scriptsize]
math0 audit: 'Math0.Statements.A195249.statement' depends on axioms [Classical.choice, Quot.sound, propext]
math0 audit: 'Math0.Frontier.A195249.result' depends on axioms [Classical.choice, Quot.sound, propext]
math0 statement check: 'Math0.Frontier.A195249.result' proves 'Math0.Statements.A195249.statement'
\end{lstlisting}

\subsection*{Toolchain and reproduction}
We use Lean~4 \cite{Lean4} \texttt{v4.33.1} and Mathlib \cite{Mathlib} \texttt{v4.33.1}, pinned in \path{lean/}; the repository also pins formal-conjectures \cite{FC} at commit \texttt{c77fdd4}, not used here.
To reproduce, run \path{scripts/verify.sh} from the repository root: it fetches the Mathlib cache, builds all math0 results with \texttt{lake build} (re-running every audit), checks that the pinned dependencies are unmodified, and reports the number of passed statement checks. The proof file alone (\texttt{lake env lean} in \path{lean}) takes under half a minute of CPU time.

{\footnotesize
\begin{thebibliography}{99}

\bibitem{DBK}
R.~Dougherty-Bliss and M.~Kauers, \emph{Hardinian arrays}, Electron.\ J. Combin.\ \textbf{31}(2) (2024), Paper No.~P2.9; arXiv:2309.00487.

\bibitem{FC}
M.~Firsching, P.~Lezeau, S.~Mercuri, M.~Z. Horv\'ath, Y.~Dillies, C.~S\"onne, E.~Wieser, F.~Zhang, T.~Hubert, B.~Ag\"uera y Arcas, and P.~Kohli, \emph{Formal Conjectures: An open and evolving benchmark for verified discovery in mathematics}, arXiv:2605.13171, 2026. Repository: \url{https://github.com/google-deepmind/formal-conjectures}.

\bibitem{Lean4}
L.~de Moura and S.~Ullrich, \emph{The Lean 4 theorem prover and programming language}, in: Automated Deduction -- CADE 28, Lecture Notes in Comput. Sci. \textbf{12699}, Springer, 2021, 625--635.

\bibitem{Mathlib}
The mathlib Community, \emph{The Lean mathematical library}, in: Proceedings of the 9th ACM SIGPLAN International Conference on Certified Programs and Proofs (CPP 2020), ACM, 2020, 367--381.

\bibitem{A195214}
The OEIS Foundation Inc., \emph{The On-Line Encyclopedia of Integer Sequences}, \url{https://oeis.org/A195214}, entry by R.~H. Hardin (September~13, 2011).

\bibitem{A195220}
The OEIS Foundation Inc., \url{https://oeis.org/A195220}, entry by R.~H. Hardin (September~13, 2011); comment by R.~Israel (October~6, 2019).

\bibitem{A195232}
The OEIS Foundation Inc., \url{https://oeis.org/A195232}, entry by R.~H. Hardin (September~13, 2011); comment by R.~Israel (November~21, 2017).

\bibitem{A195248}
The OEIS Foundation Inc., \url{https://oeis.org/A195248}, entry by R.~H. Hardin (September~13, 2011).

\bibitem{A195249}
The OEIS Foundation Inc., \url{https://oeis.org/A195249}, entry by R.~H. Hardin (September~13, 2011); conjectures by C.~Barker (May~6, 2018); revision \#9, May~6, 2018.

\end{thebibliography}
}

\end{document}
